% ==============================================================================
% Education Research Paper Template
% Structure: APA 7 + intervention design + learning-outcomes assessment
%
% APA Publication Manual (7th ed.); WhatWorks Clearinghouse Standards 4.0
% for intervention research; multilevel data (students nested in classes
% nested in schools) requires HLM.
%
% Citation style: APA 7 author-year, approximated by natbib + apalike.bst.
% Compile:  pdflatex main -> bibtex main -> pdflatex main -> pdflatex main
% ==============================================================================
\documentclass[11pt,a4paper]{article}

% ---- Page geometry ----
\usepackage[a4paper,top=25mm,bottom=25mm,left=25mm,right=25mm]{geometry}

% ---- Fonts & encoding (APA 7: Times New Roman 11pt or Georgia 11pt) ----
\usepackage[T1]{fontenc}
\usepackage[utf8]{inputenc}
\usepackage{mathptmx}
\usepackage{helvet}
\usepackage{courier}

% ---- Line numbers + double spacing for review (APA 7 manuscript) ----
\usepackage{lineno}
\usepackage{setspace}
\doublespacing
\linenumbers

% ---- Math, figures, tables ----
\usepackage{amsmath,amssymb}
\usepackage{graphicx}
\usepackage{booktabs}
\usepackage{array}
\usepackage{xcolor}
\usepackage{multirow}

% ---- Citations: APA-style author-year (approximated) ----
\usepackage[round,authoryear]{natbib}
\bibliographystyle{apalike}       % swap with plainnat or biblatex-apa for strict APA 7

% ---- Hyperref last ----
% \usepackage{hyperref}
% \hypersetup{colorlinks=true,linkcolor=blue,citecolor=blue,urlcolor=blue}

\title{A Concise and Informative Title of the Education Submission}
\author{First Author\textsuperscript{1,*}, Second Author\textsuperscript{2}, Senior Author\textsuperscript{1}\\
\textsuperscript{1}School of Education, University of Y, City, Country\\
\textsuperscript{2}Department of Curriculum \& Instruction, University of W, City, Country\\
\textsuperscript{*}e-mail: first.author@univ.edu}
\date{}

\begin{document}
\maketitle

% ==============================================================================
% Abstract  (150-250 words; APA 7)
% ==============================================================================
\begin{abstract}
% TODO: 150-250 words. Sentence 1-2: purpose and research question. Sentence
% 3-4: participants (n, grade level, setting) and design. Sentence 5-6: key
% findings with effect sizes (Cohen's d or Hedges' g) and 95\% CI. Sentence
% 7: practical implications for teachers / policy.
\end{abstract}

\noindent\textbf{Keywords:} % TODO: 3-5 keywords.
keyword one, keyword two, keyword three

% ==============================================================================
% 1. Introduction  (no "Introduction" heading per APA 7; CARS model)
% ==============================================================================
\section{Introduction}
\label{sec:intro}
% TODO: review prior literature; identify the gap; state the research
% question (\texttt{\% TODO: research-question}) and hypotheses H1, H2, ...
% Then articulate the theory of change for the intervention
% (\texttt{\% TODO: theory-of-change}).

% ==============================================================================
% 2. Method
% ==============================================================================
\section{Method}
\label{sec:method}

\subsection{Participants}
% TODO participants: describe the sampling of schools, classes, and
% students; report the nested structure ($N_{\text{schools}}$,
% $N_{\text{classes}}$, $N_{\text{students}}$); report demographic
% breakdown by treatment arm; justify sample size via power analysis.

\subsection{Research design}
% TODO: state whether the design is RCT, cluster-randomised, quasi-
% experimental (DiD, matching), or single-case; describe the
% pre-test / post-test / delayed-post-test timeline.

\subsection{Intervention}
% TODO intervention: describe the intervention in detail (manualised
% curriculum, dosage, duration, provider training); contrast with the
% business-as-usual control; reference Appendix A for the full
% intervention manual.

\subsection{Fidelity of implementation}
% TODO fidelity: describe how implementation fidelity was measured
% (observation protocol, attendance, dosage); report the fidelity
% statistics for each arm.

\subsection{Measures}
% TODO measures: list each outcome (primary: learning outcome; secondary:
% motivation, engagement) with its source, number of items, scoring,
% reliability (Cronbach's $\alpha$ in present sample), and validity
% evidence (CFA).

\subsection{Procedure}
% TODO procedure: timeline from consent -> pre-test -> intervention ->
% post-test -> delayed post-test; describe data-collection setting and
% assessor training.

\subsection{Data analysis}
% TODO analysis: specify the multilevel model (students within classes
% within schools), the fixed and random effects, the moderation and
% mediation analyses, the handling of missing data, the alpha level
% ($\alpha = 0.05$), and the effect-size metric (Cohen's $d$ / Hedges'
% $g$ with 95\% CI).

\subsection{Ethical considerations}
% TODO irb: state IRB approval number, parental informed consent
% procedure, child assent procedure, and safeguards for minors.

\subsection{Pre-registration}
% TODO preregistration: if the study was pre-registered, state the
% registry (OSF / SREE / AsPredicted) and provide the URL.

% ==============================================================================
% 3. Results
% ==============================================================================
\section{Results}
\label{sec:results}

\subsection{Descriptive statistics}
% TODO: Table 1 with means, SDs, and $n$ for each outcome by arm and time.

\begin{table}[t]
\caption{Descriptive statistics by arm and time. Caption ABOVE the table (APA convention).}
\label{tab:descriptives}
\centering
\begin{tabular}{llccc}
\toprule
Outcome & Arm & $n$ & Pre-test $M (SD)$ & Post-test $M (SD)$ \\
\midrule
Math    & Intervention & --- & --- (---) & --- (---) \\
        & Control      & --- & --- (---) & --- (---) \\
Reading & Intervention & --- & --- (---) & --- (---) \\
        & Control      & --- & --- (---) & --- (---) \\
\bottomrule
\end{tabular}
\end{table}

\subsection{Primary analysis}
% TODO: report the multilevel model estimates; reference Table 2; state
% the effect on the primary outcome with effect size and 95\% CI.

\subsection{Moderation and mediation}
% TODO: report pre-specified moderators (e.g. baseline score, gender) and
% any mediator analyses (e.g. engagement $\rightarrow$ outcome).

\subsection{Sensitivity analyses}
% TODO: re-run the primary model under alternative specifications (e.g.
% intention-to-treat vs per-protocol; alternative missing-data handling).

% ==============================================================================
% 4. Discussion
% ==============================================================================
\section{Discussion}
\label{sec:discussion}
% TODO: (i) restate the key findings without statistics, (ii) interpret in
% light of prior literature and the theory of change, (iii) theoretical
% contribution, (iv) practical implications
% (\texttt{\% TODO: practical-implications}) for teachers, school leaders,
% and policy makers, (v) limitations, (vi) future research.

\subsection{Limitations}
% TODO: sample restriction, fidelity variation, self-report measures,
% limited follow-up duration, etc.

% ==============================================================================
% 5. Conclusion
% ==============================================================================
\section{Conclusion}
\label{sec:conclusion}
% TODO: one paragraph stating the contribution and a forward-looking statement.

% ==============================================================================
% Declarations
% ==============================================================================
\section*{Declarations}

\noindent\textbf{Funding:} % TODO: grant numbers and funding bodies.

\noindent\textbf{Conflicts of interest:} % TODO: declare or state none.

\noindent\textbf{Ethics approval:} % TODO irb: IRB name + approval number.

\noindent\textbf{Informed consent:} % TODO: state parental consent + child
% assent procedure.

\noindent\textbf{Data \& materials availability:} % TODO: deposit data +
% codebook + intervention manual in a public repository (with safeguards
% for identifiable student data); state any restrictions.

\noindent\textbf{Pre-registration:} % TODO preregistration: state the
% registry URL and the date of pre-registration.

% ==============================================================================
% Acknowledgements
% ==============================================================================
\section*{Acknowledgements}
% TODO: funders, teachers, school districts, students, research assistants.

% ==============================================================================
% References  —  APA 7 author-year style
% Compile:  pdflatex main -> bibtex main -> pdflatex main -> pdflatex main
% ==============================================================================
\bibliography{references}
% Inline fallback (uncomment if not using BibTeX):
% \begin{thebibliography}{99}
% \bibitem{ref_example} Author, F., \& Author, S. (20YY). Title of the paper.
%   \emph{Journal of Educational Psychology, XX}(Y), 1--10.
% \end{thebibliography}

% ==============================================================================
% Appendices  —  intervention manual, scales, coding scheme
% ==============================================================================
\appendix

\section{Intervention Manual}
\label{app:manual}
% TODO: outline of the manualised intervention with session-by-session
% learning objectives, materials, and scripts.

\section{Measures}
\label{app:measures}
% TODO: list the full item wording for each scale and the scoring key.

\section{Coding Scheme for Process Data}
\label{app:coding}
% TODO: if process data (e.g. classroom observation) are analysed, include
% the coding scheme with definitions and exemplar codes.

\end{document}
